\section{An exact obstruction for the unmodified raw features}
The downstream transformation cannot be dropped from a deterministic
near-minimum-energy theorem. The following example attains the sharp
source energy, yet its raw frozen feature space has almost unit error
on a parity only two coordinates away. It also has polynomial-width,
arbitrarily near-minimum-energy versions.

\begin{theorem}[Raw-feature obstruction at optimal source energy]
Let $T\subset[r]$ have even cardinality $2\le j\le r-2$, with $1\in T$.
There is a width-$2^{r-1}$ network with $F=s\chi$, $E=s/\rho$, and
raw frozen feature space $\mathcal V$ satisfying
\[
 \inf_{g\in\mathcal V}\|g-\chi_T\|_2^2
                  =\frac1{1+4r b_j^2}.                 \tag{24}
\]
For $j=r-2k$ with fixed $k\ge1$, the right side tends to one and is
$1-O_k(r^{-2k-1})$.

For every $0<\varepsilon\le1$ and $0<\zeta<1$, there is also a
random construction with exactly $E=(1+\varepsilon)s/\rho$ and any width
\[
 N\ge\frac{4(1+\varepsilon)^2}{\rho^2\varepsilon^2}
                                  \log(2^r/\zeta),      \tag{25}
\]
such that, with probability at least $1-\zeta$, $\chi F\ge s$ on the
entire cube. Every realization of its raw feature space satisfies the
lower bound in (24). Thus a polynomial width and an inverse-polynomial
energy excess do not remove this obstruction.
\end{theorem}
\begin{proof}
Set $B=T\setminus\{1\}$, an odd set of size at most $r-3$, and choose
one orientation from each antipodal pair by
\[
             \mathcal Z=\{z\in\{-1,1\}^r:\chi_B(z)=1\}.
\]
Its uniform law has density $1+\chi_B$ relative to the uniform cube.
Put $h_z(x)=\relu(z\cdot x/\sqrt r)$.
For every such orientation,
\[
 \langle h_z,\chi_T\rangle=b_j\chi_T(z)=b_jz_1,
 \qquad \langle h_z,x_1\rangle=z_1/(2\sqrt r).
\]
Consequently every linear combination $g$ obeys
$\langle g,x_1\rangle=\alpha\langle g,\chi_T\rangle$,
where $\alpha=(2\sqrt r\,b_j)^{-1}$.
Orthogonality of $x_1$ and $\chi_T$ gives
\[
 \|g-\chi_T\|_2^2\ge
 (\langle g,\chi_T\rangle-1)^2+
                   \alpha^2\langle g,\chi_T\rangle^2
 \ge\alpha^2/(1+\alpha^2).
\]
The even coefficients of
$q=b_j^{-1}\E_{z\in\mathcal Z}\chi_T(z)h_z$ are exactly those
of $\chi_T$, and its only odd coefficient is $\alpha x_1$.
Indeed even characters have unchanged expectations under $1+\chi_B$;
the degree-one coefficient at $l$ is nonzero only if
$T\mathbin\triangle\{l\}=B$, equivalently $l=1$.
Thus $q=\chi_T+\alpha x_1$, and $q/(1+\alpha^2)$ attains (24).

Similarly
$\E_{z\in\mathcal Z}\chi(z)h_z=b_r\chi$: all even coefficients
except full parity vanish, and an odd coefficient would require
$[r]\mathbin\triangle\{l\}=B$, impossible because $|B|\le r-3$.
Choose $C=\sqrt{s/\rho}$, $U_z=C(0,z/\sqrt r)$, and
$a_z=C\operatorname{sgn}(b_r)\chi(z)$.
The resulting mean-field network is exactly $s\chi$, has energy $s/\rho$,
and has the same raw feature span as the $h_z$. Formula (3) gives
the asserted fixed-codimension order.

For the finite random version draw the $z_i$ independently from
$\mathcal Z$ and replace $C$ by $\sqrt{(1+\varepsilon)s/\rho}$.
Energy is deterministic, and the mean output is $(1+\varepsilon)s\chi$.
For each input the neuron output has second moment at most
$[(1+\varepsilon)s/\rho]^2$, since
$\E_{\mathcal Z}(z\cdot x)^2=r$: the conditioning character is odd
and all degree-two expectations remain zero.
Its absolute value is at most $(1+\varepsilon)s\sqrt r/\rho$.
Bernstein's inequality and a union over $2^r$ inputs show that (25)
bounds every one-sided deviation by $\varepsilon s$ with failure at
most $\zeta$. Here one uses $\rho\sqrt r\le1/2$ and
$\varepsilon\le1$ to absorb the bounded-summand term in the displayed
constant. Each sampled feature still obeys the exact linear relation,
so the obstruction holds for every realization, independently of the
margin event.
\end{proof}

The same construction also separates raw transfer from regularized
square-risk optimality. With $C^2=1/(2\rho)$ and independent orientations
from $\mathcal Z$, its energy is $1/(2\rho)$, its expected output is
$\chi/2$, and
\[
 \E_z\{\rho E+\E_x(1-\chi F)^2-3/4\}
                         \le\frac1{4\rho^2N}.
\]
Thus, for every $\eta>0$ and $N\ge1/(4\rho^2\eta)$, a realization
has population objective at most $3/4+\eta$ while retaining the same
raw-feature obstruction. This conclusion concerns the objective used in
the square-likelihood companion, without a posterior-probability claim.

This is a statement about a possible learned representation, not its
probability under the Gaussian posterior. In particular it does not
contradict the full-data, width-first raw-feature transfer theorem,
where the ordinary posterior background supplies additional features.
It shows that source interpolation, polynomial width, and almost optimal
energy alone cannot justify unmodified frozen-feature transfer.
Even projection or the linear input skip removes the explicit offending
linear component.

There is also a qualitative classification obstruction. Formula (3)
gives $|2\sqrt r\,b_j|<1$ for $r\ge6$ and the indicated $j$.
Thus the annihilating vector
$v(x)=\chi_T(x)-2\sqrt r\,b_jx_1$ has the same sign as $\chi_T$ at
every input. No function in the raw span, even with an added constant
readout bias, can classify $\chi_T$ correctly on the whole cube: such a
function would have strictly positive inner product with $v$. This
argument also covers a fixed constant tie convention. An adversary can
therefore force a mistake at every step from any deterministic learner
whose classifier is always a raw linear readout, by choosing a currently
misclassified input. This is a frozen-representation statement, not an
end-to-end Bayesian mistake lower bound.

For the full-orientation construction, the uniform population square-loss
fit has an even stronger classification failure. In the orthonormal basis
$(\chi_T,x_1)$ the raw kernel's relevant block, before its common positive
scale factor $C^2$, is
\[
 \begin{pmatrix}b_j^2&b_j/(2\sqrt r)\\
                 b_j/(2\sqrt r)&1/(4r)\end{pmatrix}.
\]
This two-dimensional subspace is invariant under the kernel operator:
the orientation law $1+\chi_B$ pairs each degree-one character with only
one even character. Consequently population ridge regression, at every
finite nonnegative regularization parameter $\gamma$, predicts a positive
multiple of
$b_j^2\chi_T+[b_j/(2\sqrt r)]x_1$.
For $\gamma=0$ this denotes the minimum-square-loss predictor in the raw
feature space. Since $|b_j|<1/(2\sqrt r)$, its classifier is exactly
$\operatorname{sgn}(b_j)x_1$, with uniform error $1/2$.
The same formula holds for a fresh Gaussian readout posterior mean after
any positive equal number of observations of every cube input, with the
corresponding positive ridge parameter. This classification conclusion
uses the explicit invariant block, not a deduction from an $L^2$ norm
lower bound alone.

Equivalently, define the full-support input law
\[
 P_T(x)=2^{-r}\{1-2\sqrt r\,b_j\chi_T(x)x_1\}.
\]
Then $\E_{P_T}\chi_Tg=0$ for every raw linear readout $g$, including
an intercept. Its population square-loss minimizer is the zero function;
its logistic loss is at least $\log2$ by Jensen. The total variation
distance between $P_T$ and uniform is $\sqrt r\,|b_j|$, which is
$\Theta_k(r^{-k-1/2})$ at fixed codimension. This is an explicit small
change of input law under which the frozen representation has no linear
target correlation.

The obstruction is stable under perturbations. Put
$v=\chi_T-2\sqrt r\,b_jx_1$, so that
$\langle h_z,v\rangle=0$ and $\|v\|_2^2=1+4r b_j^2$.
If perturbed features $\widetilde h_i$ satisfy
$N^{-1}\sum_i\|\widetilde h_i-C h_{z_i}\|_2^2\le e^2$, then
Cauchy--Schwarz gives, for every $g$ in their normalized feature RKHS,
\[
 \|g-\chi_T\|_2\le a<\|v\|_2^{-1}
 \quad\Longrightarrow\quad
 \|g\|_{\mathcal H}^2\ge
        \frac{(1-a\|v\|_2)^2}{e^2\|v\|_2^2}.           \tag{26}
\]
At $e=0$ this means that no such approximation exists.
Thus arbitrarily small perturbations restore neither a uniform bounded
readout norm nor a deterministic raw-feature analogue of (4)--(5).
